\begin{abstract}
Local field potentials across primate cortex show a spectrolaminar pattern: relative gamma power peaks superficially and relative alpha-beta power deep. Laminar cell densities predict this pattern across areas, but whether such gradients can generate it is untested. We built a three-unit cortical column of excitatory, parvalbumin (PV) and somatostatin (SOM) adaptive exponential populations. Unit sizes followed macaque laminar densities from 11 areas, connectivity scaled with density by an assumed rule, and an imposed SOM/PV ratio rose with depth. A single unit reproduced its published rhythm. Its peak frequency decreased as SOM/PV rose over 0.25--2 on seeds 1--3. E1 is inconclusive as specified: no seed oscillated at a SOM/PV ratio of 4, so the planned drop is undefined. On fresh seeds, E1b, whose band and arms were chosen after a re-read of seeds 1--3, used a 10--80 Hz peak band in two arms: arm A at fixed inhibitory count and arm B with PV fixed at 750 cells. Arm B oscillated at that ratio in all three runs, at 24--26 Hz; arm A did in one of three, at 21 Hz. Arm B passes its four criteria, arm A is inconclusive, and neither changes the E1 verdict. The limited seeds, the PV-count confound in E1 and arm A, and the co-varying total inhibition in arm B preclude attributing this shift specifically to inhibitory composition. The reference column placed alpha-beta deep on seeds 1--3 but left relative gamma near layer 4 at 0.19--0.31 of its per-frequency maximum across 12 ladder and confirmation runs, against 0.85 observed. Reference runs reached 0.23--0.26, and T8 controls 0.04--0.23. Varying network size, connectivity and SOM kinetics did not resolve this. The similarity-selected setting placed alpha-beta below the reference on selection and fresh seeds, at 0.32--0.47 against 0.83--0.86 for the reference on seeds 1--3. Flattening or reversing the density profile or the SOM/PV ramp did not shift single bands: the unit holding most relative power moved with the input, D for density changes and S for ramp changes, and the reference is also dominated by S at most frequencies. The declared three-source readout could represent the observed gamma depth profile, so this control implicates the circuit's source spectra but does not exclude observation-model effects. Moving the readout depths alone raised $r_\text{proto}$ from 0.38 at the unmoved depths, which fail the $r_{gm}$ threshold, to 0.72, the highest of 59 of 1,330 placements that met both band thresholds on the selection run; this gain is not distinguished from a frequency-shift null. The tested circuit did not reproduce the primate motif; whether architecture, parameterization or missing mechanisms cause this is unresolved. The reported tests and comparisons are uncorrected for multiple comparisons.
\end{abstract}
